\subsection{A full pro-three counterexample ramified at seven}

The implication in task B.10(b) is false, already for the actual
geometric pro-three group.  The following finite-set arithmetic
calculation is what permits us to prescribe the local cusp packet
without replacing an actual Frobenius word by an independent symbol.

\begin{statement}{B.29}{The two local factors above seven}{Proved}
Put
\[
 K=\mathbb Q(\zeta),\quad \zeta^2+\zeta+1=0,\quad
 \lambda=1-\zeta,\quad \pi=2-\zeta,\quad
 \bar\pi=3+\zeta,
\]
and let $S_K$ consist of the three primes
$(\lambda),(\pi),(\bar\pi)$.  Thus $\pi\bar\pi=7$.
Write
\[
 P_S=G_{K,S_K}^{(3)},\qquad
 L_1=G_{K_{(\pi)}}^{(3)},\qquad
 L_2=G_{K_{(\bar\pi)}}^{(3)}.
\]
After choosing prolongations, the actual continuous local maps
induce an isomorphism of pro-three groups
\[
             L_1*_3 L_2\ \xrightarrow{\ \sim\ }\ P_S.       \tag{B.29.1}
\]
Here each $L_i$ is the maximal pro-three group of $\mathbb Q_7$,
with its tame presentation
\[
 L_i=\langle t_i,F_i\mid F_it_iF_i^{-1}=t_i^7\rangle_{\rm pro\text{-}3}.
                                                               \tag{B.29.2}
\]
In particular the actual local maps are injective; this is a
conclusion, not an assumption of the construction.
Choose the two prolongations to be exchanged by a fixed complex
conjugation $c\in G_{\mathbb Q}$.  The actual quotient
\[
 \Gamma_S=
 G_{\mathbb Q,\{3,7,\infty\}}/
 \ker(G_{K,S_K}\longrightarrow P_S)
       =P_S\rtimes\langle c\rangle                             \tag{B.29.3}
\]
has $c$ exchanging the two factors of (B.29.1).
\end{statement}
\begin{proof}
Let $X=\operatorname{Spec}\mathcal O_K[1/21]$.
The Eisenstein integer ring is norm Euclidean, hence a principal
ideal domain.  Its unit group is $\mu_6$, and its localized unit
group is generated by $\mu_6,\lambda,\pi,\bar\pi$.
Kummer theory consequently identifies
\[
 H^1(P_S,\mathbb F_3)
   =H^1_{\rm et}(X,\mu_3)
   =\langle\zeta,\lambda,\pi,\bar\pi\rangle_{\mathbb F_3}.
                                                               \tag{B.29.4}
\]
We identify $\mathbb F_3$ with $\mu_3$ by $1\mapsto\zeta$.
The first equality only uses degree-one comparison: a cubic
Kummer torsor is itself a three-cover.  Each local group $L_i$
has two-dimensional $H^1$, since
$\mathbb Q_7^\times/(\mathbb Q_7^\times)^3$ has one valuation
coordinate and one residue-unit coordinate.

At $(\pi)$ the reduction of $\zeta$ is $2$, and at $(\bar\pi)$
it is $4$.  The reductions of $\lambda$ are respectively $6$
and $4$.  On arithmetic Frobenius, the Kummer character of a
unit $u$ is determined by $\bar u^{(7-1)/3}=\bar u^2$.
Thus the two Frobenius rows on $\zeta,\lambda$ are
\[
                   \begin{pmatrix}2&0\\2&2\end{pmatrix}.
                                                               \tag{B.29.5}
\]
The inertia rows on $\pi,\bar\pi$, after independently choosing
positive tame generators, form the identity matrix.  They vanish
on $\zeta,\lambda$.  Therefore the full restriction matrix has
the block form
\[
 \begin{pmatrix}
 2&0&*&*\\2&2&*&*\\0&0&1&0\\0&0&0&1
 \end{pmatrix},\qquad \det=1\quad\text{in }\mathbb F_3.          \tag{B.29.6}
\]
The unspecified entries retain the actual choices of Frobenius;
they have no effect on this determinant.  By the universal
property of the free pro-three product, (B.29.6) says that
restriction from $P_S$ to $B=L_1*_3L_2$ is an isomorphism in
$H^1(-,\mathbb F_3)$.

We also have an injection
\[
 H^2(P_S,\mathbb F_3)\lhook\joinrel\longrightarrow
 H^2_{\rm et}(X,\mu_3).                                      \tag{B.29.7}
\]
This is the general degree-two injection preceding
\citet[Proposition 2.1, author version p.~4]{CBSchmidt2007}:
the first cohomology of the universal pro-three cover vanishes,
and the Hochschild--Serre five-term sequence gives the injection.
For clarity, a cyclic three-torsor on a finite Galois three-cover
is killed on a further Galois three-cover: take its finitely
many conjugates and their compositum.  Its Galois group over
the original base is an extension of a subgroup of a product
of cyclic groups of order three by a three-group.  This proves
the vanishing used in that five-term argument.  We do not need
a blanket comparison in all cohomological degrees.

Since $\operatorname{Pic}(X)=0$, the Kummer sequence gives
$H^2_{\rm et}(X,\mu_3)=\operatorname{Br}(X)[3]$.
The Brauer group of the regular integral curve $X$ injects
into $\operatorname{Br}(K)$
\citep[Theorem 3.5.4]{CBColliotTheleneSkorobogatov2019}.
An element of this Brauer group has zero local invariant outside
$S_K$: at a finite point of $X$ it extends over the completed
local integer ring, whose Brauer group is zero because that
ring is henselian with finite residue field
\citep[Theorem 3.4.2]{CBColliotTheleneSkorobogatov2019}.
There are no real places.  Global Brauer reciprocity
\citep[VIII, Theorem 4.2]{BMilneCFT2020} therefore implies that
restriction to the two primes above seven is injective on
$\operatorname{Br}(X)[3]$.  Indeed, if those two invariants
vanish, the sole remaining invariant at $(\lambda)$ is zero
by their sum, and all local invariants then vanish.
By naturality of (B.29.7), a class in $H^2(P_S,\mathbb F_3)$
whose restrictions to $L_1$ and $L_2$ are zero is zero.
In particular
\[
                H^2(P_S,\mathbb F_3)\longrightarrow
                       H^2(B,\mathbb F_3)
                      \quad\text{is injective}.              \tag{B.29.8}
\]
This argument does not assume in advance that either local
group embeds in $P_S$.

Let $f:B\to P_S$ be the actual local map.  Its image is compact
and hence closed.  If this image were proper, its image would be
proper in a finite three-group quotient of $P_S$ and would lie
in a maximal subgroup of index three.  The corresponding
nonzero $H^1(P_S,\mathbb F_3)$ character would restrict to zero
on $B$, contradicting (B.29.6).  Thus $f$ is surjective.
For $R=\ker f$, its five-term sequence and (B.29.6)--(B.29.8)
give
\[
                  H^1(R,\mathbb F_3)^{B}=0.                  \tag{B.29.9}
\]
If $R$ were nontrivial, choose a finite three-group quotient
$B_0$ of $B$ in which its image $R_0$ is nontrivial.
The nonzero module $R_0/\Phi(R_0)$ has nonzero coinvariants
under $B_0$: the augmentation ideal of $\mathbb F_3[B_0]$ is
nilpotent.  A nonzero linear functional on those coinvariants
gives a nonzero $B$-invariant character of $R$, contrary to
(B.29.9).  This proves (B.29.1).

For (B.29.2), wild inertia is pro-seven and hence dies in every
pro-three quotient.  Tame inertia has three-primary part
$\mathbb Z_3(1)$ and the unramified quotient has three-primary
part $\mathbb Z_3$.  Since $7\equiv1\pmod3$, its action by
$7$ is continuous on this quotient; the tame presentation
therefore gives exactly (B.29.2).
Finally $K/\mathbb Q$ is Galois and is ramified only at three,
so $G_{K,S_K}$ is the actual index-two subgroup of the rational
restricted-ramification group in (B.29.3).
Its maximal-pro-three kernel is characteristic and hence normal
in that group.  The image of the chosen complex conjugation
has order two and splits the quotient.  Taking conjugate
prolongations makes its exchange of local factors literal.
There is no independent prescription of a higher Frobenius
word anywhere in the argument.
\end{proof}

\begin{statement}{B.30}{Ramification at seven does not force a rational cusp}{Refuted}
Let
\[
 U_3=\{X^3+Y^3+Z^3=0\}\setminus\{XYZ=0\},\qquad
 N=\pi_1(U_{3,\overline{\mathbb Q}})^{(3)}.
\]
Use the genuine arithmetic action based at the rational unit
tangent at $O=[0:-1:1]$.  The assertion that every section
ramified at a prime different from three belongs to a global
rational-cusp packet is false.  A counterexample is supplied
by each nonzero member of the following family.
There is a family, indexed by
$a\in\mathbb Z_3$, of actual continuous sections
\[
                    s_a:G_{\mathbb Q}\longrightarrow
                                  N\rtimes G_{\mathbb Q}
\]
with pairwise distinct geometric conjugacy classes, unramified
outside $\{3,7\}$.  At seven, $s_a$ is the actual local cusp
section centred at
\[
                         P=[0:-\zeta:1]
\]
in the embedding $\zeta\equiv2\pmod7$, with tangent twisted by
$a\kappa_3(7)$.  For every $a\ne0$ it is ramified at seven
and is not the image of a global rational cusp section.

The same conclusions hold after projecting to every quotient
$N/J$ for which $J$ is closed, normal, $G_{\mathbb Q}$-stable,
and contained in $N'$.  In particular task B.10(b), and the
corresponding rational-cusp implication in B.10(c) for
$N/N''$ and the marked quotient, are false.  This example has
only one ramified rational prime away from three; it makes
no assertion about equality of colours at distinct rational
primes or about infinitely many ramified primes.
\end{statement}
\begin{proof}
We first specify the based action, because using merely an
outer action would not suffice for this construction.
The cubic power cover of the tripod identifies the actual
geometric group with the characteristic open subgroup
\[
        N=\ker\bigl(F_2^{(3)}\longrightarrow(\mathbb Z/3)^2\bigr).
                                                               \tag{B.30.1}
\]
Its rank is $1+9(2-1)=10$.
At $O$ the lift of the standard tripod tangent is given by
$X/Z=-s$, $t=s^3$, $Y/Z=-(1-s^3)^{1/3}$.
Coefficientwise Galois action on the coherent Puiseux fibres
makes (B.30.1) an identification for the genuine based
arithmetic action, as proved in B.27.
Over $K$ the action on the Frattini quotient of $F_2^{(3)}$
is trivial: it is the mod-three cyclotomic scalar.
The kernel of the action on that Frattini quotient is pro-three,
by the finite three-group basis theorem applied to successive
characteristic finite quotients.
Its compact arithmetic subgroup restricts continuously to $N$.
The pair and the chosen unit tangent have good reduction
outside three, so inertia there acts trivially on $N$ by
prime-to-residue-characteristic specialization
\citep[XIII, Corollary 2.12]{BSGA1}; this is the based
specialization argument of B.27.
Consequently the actual action $\rho_N$ factors through
$\Gamma_S$ of (B.29.3), and its restriction to $P_S$ has
pro-three image.  This is a nonlinear statement about the
actual action on the entire free pro-three group, not only
about its abelianization.

The non-rational cusp $P$ is $K$-rational, and in the chosen
completion it is $\mathbb Q_7$-rational.  Choose its unit
tangent using the local parameter $X/Z$.
In the geometric pro-three fundamental exact sequence over
$\mathbb Q_7$, choose a geometric path from the fixed tangent
at $O$ to this tangent at $P$.  Its tangential section,
twisted inside its peripheral group by $a\kappa_3(7)$,
defines an actual continuous cocycle
\[
                z_{1,a}:G_{\mathbb Q_7}\longrightarrow N
\]
for the restriction of the original action $\rho_N$.
All paths are kept fixed as $a$ varies.  The local action has
pro-three image, so the graph homomorphism
\[
 g\longmapsto(z_{1,a}(g),\rho_N(g))
       \quad\text{in }N\rtimes\rho_N(G_{\mathbb Q_7})
                                                               \tag{B.30.2}
\]
has a pro-three target.  It factors through the actual maximal
pro-three quotient $L_1$.  We retain the same notation for its
cocycle there.  Under the literal local embedding in B.29,
$h\mapsto(z_{1,a}(h),h)$ is a continuous homomorphism
$L_1\to N\rtimes P_S$ whose second projection is the inclusion.

Choose the second local cocycle, on the conjugate factor,
by the exact formula
\[
                 z_{2,a}(chc)=\rho_N(c)z_{1,a}(h)
                         \qquad(h\in L_1).                    \tag{B.30.3}
\]
It is the actual local cusp section at $\bar P=cP$,
with the conjugate path and tangent.  The cocycle identity
follows from
$\rho_N(chc)=\rho_N(c)\rho_N(h)\rho_N(c)$.
We do not identify a positively normalized tame generator at
the second prime with $ct_1c$; complex conjugation acts on
the Tate twist and that normalization is unnecessary.

Both local graph sections now have target $N\rtimes P_S$,
a pro-three group.  By (B.29.1) and the universal property of
the free pro-three product they extend uniquely to
\[
            \sigma_a:P_S\longrightarrow N\rtimes P_S.
                                                               \tag{B.30.4}
\]
Its arithmetic projection is the identity, because it is
the identity on both generating local factors.  Write
$\sigma_a(h)=(z_a(h),h)$.  Formula (B.30.3) and uniqueness
give $z_a(chc)=\rho_N(c)z_a(h)$ on all of $P_S$.
Therefore setting $z_a(c)=1$ extends this to a continuous
cocycle on $\Gamma_S=P_S\rtimes\langle c\rangle$.
Let $r:G_{\mathbb Q}\to\Gamma_S$ be the actual quotient map.
The required section is
\[
                  s_a(g)=(z_a(r(g)),g).                       \tag{B.30.5}
\]
The second coordinate here is $g$, not $r(g)$.
Thus (B.30.5) is a genuine section of the original arithmetic
extension.  It is the action and the geometric cocycle which
factor through $\Gamma_S$, not the section's identity
arithmetic projection.  At every prime outside $\{3,7\}$
both the action and this cocycle are trivial on inertia.
Its localization at seven is the chosen local packet.

We finish with a conjugacy-invariant obstruction to global
cuspidality, already in $V=N^{\rm ab}$.
Let $D$ be the nine geometric coordinate cusps.  The usual
homology sequence of the punctured smooth cubic gives an
injection
\[
 W=\mathbb Z_3(1)[D]\big/
          \mathbb Z_3(1)\sum_{d\in D}e_d
           \lhook\joinrel\longrightarrow V.                  \tag{B.30.6}
\]
The generator of inertia about a cusp $d$ has image $e_d$.
At seven the unit tangential section at $P$ is unramified;
this follows from the smooth integral pair and its integral
unit tangent.  Choose $t_1$ so that $\kappa_3(7)(t_1)=1$.
Then the restriction of the abelian cocycle of $s_a$ to
this inertia generator is exactly
\[
                               a e_P.                        \tag{B.30.7}
\]
The geometric path only conjugates peripheral generators,
which has no effect in $V$.
Since inertia at seven acts trivially on $V$, a geometric
coboundary has zero restriction there.  Hence (B.30.7) is
invariant under global or local geometric conjugacy.
It also distinguishes all parameters $a$, since $e_P$
has infinite order.

The rational coordinate cusps are precisely
$[0:-1:1],[-1:0:1],[-1:1:0]$.
For a global rational cusp packet centred at such a cusp
$C$, its inertia value at seven belongs to
$\mathbb Z_3e_C$; changing a tangent merely changes that
scalar.  If $P\ne C$, the two classes $e_P,e_C$ are linearly
independent in (B.30.6): compare a third coordinate in
any alleged relation with the diagonal vector, then compare
the $P$ and $C$ coordinates.  Thus for $a\ne0$ equation
$ae_P=be_C$ is impossible for every $b\in\mathbb Z_3$.
This proves noncuspidality.

If $J\subseteq N'$, the quotient $N/J$ retains the entire
abelianization $V$.  The preceding obstruction therefore
survives its projection, as do ramification and distinctness
of parameters.  The intrinsic metabelian kernel $N''$ is
contained in $N'$.  In the exact marked notation of
LM7 E.71, $I_Q$ is the closed normal closure of a peripheral
generator $d_Q$ at the rational marked cusp $Q$, and
$J=\overline{\gamma_3(I_Q)}$.
This is normal in $N$ and Galois-stable: $I_Q$ is normal
and stable, and its closed lower central subgroups are
characteristic in it.  It is contained in $N'$, since its
defining triple commutators are.  Thus the same argument
applies to that specified marked quotient.
This transfer uses no assertion that
all sections of a quotient lift.
\end{proof}

\paragraph{Exact scope of the counterexample.}
The program \path{c_seven_free_product.py} independently
checks the Eisenstein norm identities, all residue-symbol
entries of (B.29.5), all $81$ possible upper-right blocks
in (B.29.6), and primitive independence of the relevant
boundary vectors.  The passage from these finite checks
to the full pro-three group is the cohomological proof
of B.29, not a computational inference.
If the chosen cusp in (B.30.2) is rational instead and the
local path is the restriction of a representative of its
global packet with value one at $c$, this conjugate gluing
recovers that global packet twisted by $a\kappa_3(7)$.
Indeed that global packet then has exactly the prescribed
restrictions to both local factors and the chosen value
at $c$, so uniqueness in (B.30.4) applies.
Without this compatibility on paths, local conjugacy is
not equality of the prescribed local graph maps, and no
global-cuspidality conclusion is inferred from uniqueness.
In either case this construction has only one possible
ramified rational prime away from three, so it does not
test equality of colours at different rational primes.
Tasks B.10(a) and the full pro-three form of B.11 remain
open in this return.  Adding an auxiliary rational prime
to the present construction introduces its actual
Frobenius word; the four-generator free-product proof
does not make that extra word independent.
